% Part: history
% Chapter: biographies 
% Section: emmy-noether

\documentclass[../../../include/open-logic-section]{subfiles}

\begin{document}

\olfileid{his}{bio}{noe}

\olsection{ఎమ్మీ నోయెథర్}

\olphoto{noether-emmy}{ఎమ్మీ నోయెథర్}

ఎమ్మీ నోయెథర్ (ఉచ్చారణ: నెర్-టర్) 1882 మార్చి 23న జర్మనీలోని ఎర్లాంగెన్‌లో, విద్యావంతులైన ఎగువ మధ్యతరగతి కుటుంబంలో జన్మించారు. ``ఆధునిక బీజగణితానికి తల్లి''గా ప్రశంసలు పొందిన నోయెథర్, మహిళల విద్యకు గణనీయమైన అడ్డంకులు ఉన్నప్పటికీ గణితం, భౌతికశాస్త్రం రెండింటికీ మార్గదర్శకమైన కృషి చేశారు. అప్పటి జర్మనీలో బాలికలకు కళల్లో విద్య ఇవ్వాలని భావించేవారు; కళాశాల ప్రవేశానికి సిద్ధం చేసే పాఠశాలల్లో వారికి అనుమతి ఉండేది కాదు. అయితే గ్యోటింగెన్, ఎర్లాంగెన్ విశ్వవిద్యాలయాల్లో (రెండవదానిలో ఆమె తండ్రి గణిత ఆచార్యుడు) తరగతులకు శ్రోతగా హాజరైన తరువాత, మహిళా విద్యార్థులను అనుమతించేలా విధానం మారిన 1904లో నోయెథర్ ఎర్లాంగెన్‌లో విద్యార్థినిగా చేరగలిగారు. 1907లో గణితంలో డాక్టరేట్ పొందారు.

అర్హతలున్నప్పటికీ నోయెథర్ తన వృత్తి జీవితంలో తీవ్ర వ్యతిరేకతను ఎదుర్కొన్నారు. 1908--1915 మధ్య ఎర్లాంగెన్‌లో జీతం లేకుండా బోధించారు. ఆ సమయంలోనే అప్పటి ప్రపంచ ప్రముఖ గణితశాస్త్రజ్ఞులలో ఒకరైన డేవిడ్ హిల్బర్ట్ దృష్టిని ఆకర్షించారు; ఆయన ఆమెను గ్యోటింగెన్‌కు ఆహ్వానించారు. అయితే మహిళలకు ఆచార్య పదవులు నిషేధించబడటంతో, ఆమె హిల్బర్ట్ పేరుతోనే, మళ్లీ జీతం లేకుండానే, ఉపన్యాసాలివ్వగలిగారు. ఈ కాలంలోనే నేడు నోయెథర్ సిద్ధాంతం అని పిలిచే, ఇప్పటికీ సిద్ధాంత భౌతికశాస్త్రంలో వాడే ఫలితాన్ని నిరూపించారు. 1919లో చివరకు నోయెథర్‌కు బోధించే హక్కు లభించింది. విశ్వవిద్యాలయ సహచరుల నిరంతర వ్యతిరేకతకు హిల్బర్ట్ ``మహాశయులారా, అధ్యాపక మండలి స్నానశాల కాదు'' అని ప్రతిస్పందించినట్లు చెబుతారు.

1920ల చివర్లో ఆమె అమూర్త బీజగణితంపై దృష్టి పెట్టారు; ఆమె కృషి ఆ రంగంలో విప్లవాత్మక మార్పు తెచ్చింది. నిరూపణలలో ఆమె తరచుగా ఆరోహణ శ్రేణి షరతు అని పిలిచేదాన్ని వాడారు; కొన్ని సమితుల కచ్చితంగా పెరుగుతూ పోయే అనంత శ్రేణి ఉండదని అది చెబుతుంది. ఉదాహరణకు, నోయెథరియన్ వలయాలు అని ఇప్పుడు పిలిచే కొన్ని బీజగణిత నిర్మాణాల్లో ఐడియళ్ల $I_1
\subsetneq I_2 \subsetneq \dots$ వంటి అనంత శ్రేణులు ఉండవు. ఈ షరతును ఏ పాక్షిక క్రమానికైనా సాధారణీకరించవచ్చు (బీజగణితంలో ఇది ఉపసమితి సంబంధంతో క్రమబద్ధమైన ఐడియళ్ల ప్రత్యేక సందర్భం); అలాగే పాక్షిక క్రమంలో కచ్చితంగా \emph{తగ్గుతూ} పోయే ప్రతి శ్రేణీ చివరకు ఆగిపోయే ద్వైత అవరోహణ శ్రేణి షరతును కూడా పరిశీలించవచ్చు. ఒక పాక్షిక క్రమం అవరోహణ శ్రేణి షరతును తృప్తిపరస్తే,~$\Nat$పై $<$ క్రమం వెంట ఆగమనం వాడినట్లే ఆ క్రమం వెంట కూడా ఆగమనాన్ని వాడవచ్చు. ఇటువంటి క్రమాలను \emph{సుస్థాపిత} లేదా \emph{నోయెథరియన్} అంటారు; అనురూప నిరూపణ సూత్రాన్ని \emph{నోయెథరియన్ ఆగమనం} అంటారు.

నోయెథర్ యూదు మతస్థురాలు. 1933లో నాజీలు అధికారంలోకి వచ్చినప్పుడు ఆమెను పదవి నుంచి తొలగించారు. అదృష్టవశాత్తూ, అమెరికా సంయుక్త రాష్ట్రాల్లో పెన్సిల్వేనియాలోని బ్రిన్ మావర్‌లో తాత్కాలిక పదవికి నోయెథర్ వలస వెళ్లగలిగారు. అక్కడ ఉన్నప్పుడు ప్రిన్స్‌టన్‌లో కూడా ఉపన్యాసాలిచ్చారు; అయితే ఆ విశ్వవిద్యాలయం మహిళలకు ఆహ్వానకరంగా లేదని భావించారు \citep[81]{Dick1981}. 1935లో గర్భాశయ కణితిని తొలగించడానికి నోయెథర్‌కు శస్త్రచికిత్స జరిగింది. ఆ శస్త్రచికిత్స కారణంగా వచ్చిన సంక్రమణతో ఆమె మరణించారు; బ్రిన్ మావర్‌లోనే సమాధి చేశారు.

\begin{reading}
నోయెథర్ జీవిత చరిత్రకు \citet{Dick1981} చూడండి. పెరిమీటర్ ఇన్‌స్టిట్యూట్ ఫర్ థియరెటికల్ ఫిజిక్స్, ఆమె జీవితం, ప్రభావం గురించి ఉపన్యాసాలను ఆన్‌లైన్‌లో అందుబాటులో ఉంచింది \citep{Perimeter2015}. చదవడం నుంచి విరామం కావాలంటే, \emph{Stuff You Missed in History Class}లో నోయెథర్ జీవితం, ప్రభావంపై పాడ్‌కాస్ట్ ఉంది \citep{Frey2015}. నోయెథర్ సంకలిత రచనలు మూల జర్మన్‌లో లభిస్తాయి \citep{Noether1983}.
\end{reading}

\end{document}
